% Generated from locale/tr/content/first-order-logic/proof-systems/introduction.tex; do not hand-edit.


\iftag{FOL}
      {\olfileid[tr]{fol}{prf}{int}}
      {\olfileid[tr]{pl}{prf}{int}}

\olsection{Giriş}

Mantıkların genellikle hem bir semantiği hem de !!a{derivation} sistemi
vardır. Semantik; doğruluk, sağlanabilirlik, geçerlilik ve mantıksal gerektirme
gibi kavramlarla ilgilenir. !!{derivation} sistemlerinin amacı, mantıksal
gerektirmeyi ve geçerliliği belirlemenin salt sözdizimsel bir yöntemini
sağlamaktır. Bu sistemler salt sözdizimseldir; çünkü böyle bir sistemdeki
!!{derivation}, genellikle !!{sentence}s ya da !!{formula}s içeren bir dizi
(veya başka bir sonlu düzenleme) biçiminde, sonlu bir sözdizimsel nesnedir.
İyi !!{derivation} sistemlerinde, !!{sentence}s ya da !!{formula}s içeren
verilmiş herhangi bir dizi veya düzenlemenin “doğru” olup olmadığı mekanik
olarak denetlenebilir.

\iftag{FOL}{Birinci dereceden mantığın}{Önerme mantığının} en basit
(ve tarihsel olarak ilk) !!{derivation}
sistemleri \emph{aksiyomatik}ti. Böyle bir sistemde !!{formula}s içeren bir
dizi şu koşullarda !!a{derivation} sayılır: dizideki her !!{formula} ya sabit
bir “aksiyom” kümesinde yer alır ya da daha önce gelen !!{formula}s içinden
sabit sayıdaki “çıkarım kurallarından” biriyle elde edilir. Ayrıca !!a{formula} aksiyom
mudur ve diğer !!{formula}s üzerinden bir çıkarım kuralıyla doğru biçimde
elde edilmiş midir, bunlar mekanik olarak denetlenebilir. Aksiyomatik
!!{derivation} sistemlerini betimlemek ve metakuramsal olarak ele almak
kolaydır; ancak bu sistemlerde !!{derivation}s üretmek de, onları okuyup
anlamak da zordur.

Başka !!{derivation} sistemleri, çeşitli !!{derivation}s kurmayı ya da
!!{derivation}s tamamlandıktan sonra onları anlamayı kolaylaştırmak amacıyla
geliştirilmiştir. Doğal tümdengelim, tableau ispatları olarak da bilinen
doğruluk ağaçları ve sekant hesabı bunlara örnektir. Bazı !!{derivation}
sistemleri özellikle mekanikleştirme gözetilerek tasarlanır; örneğin
çözümleme yöntemini yazılımda uygulamak kolaydır (ama ortaya çıkan
!!{derivation}s neredeyse anlaşılamaz). Bu diğer !!{derivation}
sistemlerinin çoğu, !!{derivation}s için diziler yerine !!{formula}s içeren
ağaçlar kullanır. Böylece !!a{derivation} içindeki hangi parçaların hangi
başka parçalara dayandığını görmek kolaylaşır.

Dolayısıyla \iftag{FOL}{birinci dereceden mantık}{önerme mantığı} gibi
belirli bir mantık için farklı
!!{derivation} sistemleri, !!a{sentence} ne zaman \emph{teorem} sayılacağına
ve !!a{sentence} başka tümcelerden ne zaman !!{derivable} olduğuna farklı
açıklamalar getirir. Bu nasıl yapılırsa yapılsın (aksiyomatik !!{derivation}s,
doğal tümdengelimler, sekant !!{derivation}s, doğruluk ağaçları ya da
çözümlemeyle çürütmeler aracılığıyla), bu bağıntıların semantik geçerlilik ve
mantıksal gerektirme kavramlarıyla örtüşmesini isteriz. “$!A$ bir teoremdir”
için $\Proves !A$, “$!A$'nın $\Gamma$'dan !!{derivable} olduğu” içinse
“$\Gamma \Proves !A$” yazalım. $\Proves$ nasıl tanımlanırsa tanımlansın,
$\Entails$ ile örtüşmesini, yani şunların geçerli olmasını isteriz:
\begin{enumerate}
\item $\Proves !A$ ancak ve ancak $\Entails !A$
\item $\Gamma \Proves !A$ ancak ve ancak $\Gamma \Entails !A$
\end{enumerate}
Yukarıdaki ifadelerin “yalnızca eğer” yönüne \emph{sağlamlık} denir.
Bir !!^a{derivation} sistemi, !!{derivability} mantıksal gerektirmeyi (veya
geçerliliği) güvence altına alıyorsa sağlamdır. Her makul !!{derivation}
sistemi sağlam olmalıdır; sağlam olmayan !!{derivation} sistemleri hiçbir işe
yaramaz. Sonuçta !!a{derivation} vermenin bütün amacı, geçerlilik veya
mantıksal gerektirme için sözdizimsel bir güvence sağlamaktır. Sunduğumuz
!!{derivation} sistemlerinin sağlamlığını ispatlayacağız.

Bunun tersi olan “eğer” yönü de önemlidir; buna \emph{tamlık} denir. Tam bir
!!{derivation} sistemi, $!A$ geçerli olduğunda $!A$'nın teorem olduğunu ve
$\Gamma \Entails !A$ olduğunda $\Gamma \Proves !A$ olduğunu gösterecek kadar
güçlüdür. Tamlığı ortaya koymak daha zordur ve bazı mantıklarda tam bir
!!{derivation} sistemi yoktur.
\iftag{FOL}{Birinci dereceden mantıkta vardır. Kurt G\"odel, 1929 tarihli
doktora tezinde birinci dereceden mantığın !!a{derivation} sistemi için
tamlığı ispatlayan ilk kişiydi.}{Önerme mantığında tam türetim sistemleri
vardır.}

!!{derivation} sistemleriyle bağlantılı başka bir kavram da
\emph{tutarlılık}tır. !!{sentence}s içeren bir kümeden her şeyi
!!{derive} mümkünse bu kümeye tutarsız, aksi halde tutarlı denir. Tutarsızlık,
sağlanamazlığın sözdizimsel karşılığıdır: sağlanamaz kümeler gibi,
!!{sentence}s içeren tutarsız kümeler de iyi kuramlar oluşturmaz; temel bir
kusurları vardır. !!{sentence}s içeren tutarlı kümeler doğru veya yararlı
olmayabilir, ama en azından mantıksal yararlılığın bu asgari eşiğini aşarlar.
Farklı !!{derivation} sistemlerinde, !!{sentence}s içeren kümelerin özel
tutarlılık tanımı değişebilir; yine de $\Proves$ için olduğu gibi,
tutarlılığın semantik karşılığı olan sağlanabilirlikle örtüşmesini isteriz.
$\Gamma$'nın ancak ve ancak sağlanabilir olduğunda tutarlı olmasını isteriz.
Burada “yalnızca eğer” yönü tamlığa (tutarlılık sağlanabilirliği güvence altına
alır), “eğer” yönü ise sağlamlığa (sağlanabilirlik tutarlılığı güvence altına
alır) karşılık gelir. Nitekim klasik
\iftag{FOL}{birinci dereceden mantıkta}{önerme mantığında} sağlamlık ve
tamlığın bu iki biçimi birbirine denktir.

